\documentclass[10pt,a4paper]{amsart}
\usepackage[margin=19mm]{geometry}
\usepackage{amsmath,amssymb,mathtools}
\usepackage[scr=rsfs,cal=euler]{mathalfa}
\usepackage{xfrac}
\usepackage[unicode,hidelinks,bookmarks=false]{hyperref}
\newcommand{\ie}{i.e.\ }
\newcommand{\cf}{cf.\ }
\newcommand{\resp}{respectively\ }
\newcommand{\Subsection}{\subsection}
\providecommand{\nobreakdash}{\nobreak\hbox{-}}
\setlength{\emergencystretch}{2em}
\multlinegap0pt
\usepackage{amssymb}
\newtheorem{lemma}{Lemma}
\newtheorem{proposition}{Proposition}
\title{On two Romanoff type problems of Erd\H{o}s}
\author{Yuchen Ding}
\date{Source: arXiv:2503.22700v5 (CC BY 4.0)}
\begin{document}
\maketitle

\noindent\emph{Source and licence.} On two Romanoff type problems of Erd\H{o}s. Version arXiv:2503.22700v5 (CC BY 4.0), distributed by arXiv under the Creative Commons Attribution 4.0 International licence, http://creativecommons.org/licenses/by/4.0/, as recorded in the arXiv licence metadata for this exact version and shown on the abstract page https://arxiv.org/abs/2503.22700v5. Full licence notice: \texttt{licenses/arxiv-2503.22700-CC-BY-4.0.txt}.\par\smallskip

\noindent\textit{Editorial scope.} Counted unit: the article's proposition prop-metric (source0.tex lines 591--600). The article's complete original proof of it is delivered with it (source0.tex lines 602--805, one \texttt{proof} environment of 204 lines). No mathematics is written, repaired or re-derived: the statement and the proof are the article's own lines, in the article's own notation, and no retained line is substituted or re-flowed.  Retained uncounted as the source-local context the proof needs: lines 282--288; lines 292--299; lines 409--415; lines 470--482; lines 510--517.  Nothing was omitted from any retained span, so the package carries no omission record.  No retained line contains a citation, so the wrapper carries no bibliography and no bibliography entry.  No retained line contains a figure environment, an image-inclusion command or drawing code, so no image is delivered and the package declares no asset dependency.  Editorial formatting only, in addition: an AMS \texttt{amsart} preamble that restates the article's own declarations and macros used by the retained lines, namely the environments \texttt{lemma}, \texttt{proposition} on separate counters, together with the standard packages those lines need.  The retained environments print in this document as Lemma 1, Proposition 1 in order of appearance, whereas the article prints the same items with the numbering of its own full document.  The complete native source of the article as distributed in the arXiv e-print for this exact version is bundled unchanged as \texttt{sources/175578/source0.tex}.
    \begin{lemma}\label{lem-borel-cantelli}
        Let $\mathcal{E}_1,\mathcal{E}_2,\ldots$ be measurable sets. If
        $$
        \sum_{m\ge1}\operatorname{meas}(\mathcal{E}_m)<\infty,
        $$
        then almost every point belongs to only finitely many of the sets $\mathcal{E}_m$.
    \end{lemma}

    \begin{lemma}\label{lem-erdos-turan}
        Let $x_1,\ldots,x_K$ be real numbers and let $I\subset[0,1)$ be an interval. For every $H\ge1$,
        $$
        \left|\#\bigl\{1\le k\le K\mid \{x_k\}\in I\bigr\}-K|I|\right|
        \le C\left(\frac{K}{H}+\sum_{1\le h\le H}\frac1h\left|\sum_{k\le K}e(hx_k)\right|\right),
        $$
        where $C$ is an absolute constant.
    \end{lemma}

    \begin{lemma}\label{lem-exp-sum}
        Let $J=[Y_1,Y_2]$, where $1<Y_1<Y_2$, and let $B>0$ be fixed. If $K\ge2$, $1\le d\le K^B$, and $1\le h\le K$, then
        $$
        \int_J\left|\sum_{k\le K}e\left(\frac{h y^k}{d}\right)\right|^2\,\mathrm{d}y
        \ll_{Y_1,Y_2,B} K.
        $$
    \end{lemma}

    \begin{lemma}\label{lem-collision}
        Let $J=[Y_1,Y_2]$, where $1<Y_1<Y_2$. For $1\le k<\ell$ and $d\ge1$, define
        $$
        \mathcal{E}_{k,\ell}(d)=\bigl\{y\in J\mid
        \lfloor y^\ell\rfloor\equiv \lfloor y^k\rfloor\pmod d\bigr\}.
        $$
        Then
        $$
        \operatorname{meas}(\mathcal{E}_{k,\ell}(d))
        \ll_{Y_1,Y_2}
        \frac1d+\frac1{(\ell-k)Y_1^{\ell-1}}.
        $$
    \end{lemma}

    \begin{lemma}\label{lem-residue-ud}
        For almost every $y>1$, for every integer $q\ge1$ and every $0\le r<q$,
        $$
        \#\{1\le k\le K\mid \lfloor y^k\rfloor\equiv r\pmod q\}
        =\frac Kq+o_{y,q}(K)
        $$
        as $K\to\infty$.
    \end{lemma}

    \begin{proposition}\label{prop-metric}
        Let $J=[Y_1,Y_2]$, where $1<Y_1<Y_2$. For almost all $y\in J$,
        \begin{equation}\label{eq-pair-asymptotic}
            \sum_{\substack{1\le k<\ell\le K\\
                    \lfloor y^k\rfloor\ne\lfloor y^\ell\rfloor}}
            W\!\left(\lfloor y^\ell\rfloor-\lfloor y^k\rfloor\right)
            =\left(\frac{\zeta(2)}{2\zeta(4)}+o_y(1)\right)K^2
        \end{equation}
        as $K\to\infty$.
    \end{proposition}

    \begin{proof}
        Choose $B>2$ so large that
        \begin{equation}\label{eq-B-choice}
            B\frac{\log Y_1}{\log Y_2}>3.
        \end{equation}
        The implied constants below may depend on $J$ and $B$.

        We begin with squarefree moduli not exceeding a power of the number of terms. Fix $\lambda>1$ and put
        $$
        L_m=\lfloor \lambda^m\rfloor,\qquad D_m=L_m^B.
        $$
        For $d\le D_m$, the congruence $\lfloor y^k\rfloor\equiv a_0\pmod d$ is equivalent to
        $$
        \left\{\frac{y^k}{d}\right\}\in
        \left[\frac{a_0}{d},\frac{a_0+1}{d}\right)
        $$
        for the representative $a_0\in\{0,1,\ldots,d-1\}$. Applying Lemma~\ref{lem-erdos-turan} with $H=L_m$ gives
        $$
        M_{L_m,d}(y)
        \le \frac{L_m}{d}
        +C
        +C\sum_{1\le h\le L_m}\frac1h
        \left|\sum_{k\le L_m}e\left(\frac{h y^k}{d}\right)\right|.
        $$
        Consequently
        \begin{equation}\label{eq-small-moduli-sharp}
            \sum_{\substack{d\le D_m\\ d\ \mathrm{squarefree}}}\frac{M_{L_m,d}(y)}d
            \le L_m\sum_{\substack{d\le D_m\\ d\ \mathrm{squarefree}}}\frac1{d^2}+O(\log L_m)+T_m(y),
        \end{equation}
        where
        $$
        T_m(y)=\sum_{\substack{d\le D_m\\ d\ \mathrm{squarefree}}}\sum_{h\le L_m}\frac1{dh}
        \left|\sum_{k\le L_m}e\left(\frac{h y^k}{d}\right)\right|.
        $$
        By Cauchy's inequality and Lemma~\ref{lem-exp-sum},
        \begin{align*}
            \int_J T_m(y)\,\mathrm{d}y\!
            \le
            \sum_{\substack{d\le D_m\\ d\ \mathrm{squarefree}}}\sum_{h\le L_m}\frac{1}{dh}
            \left(\operatorname{meas}(J)\int_J\left|\sum_{k\le L_m}e\left(\frac{h y^k}{d}\right)\right|^2\,\mathrm{d}y\right)^{1/2} 
            \!\ll L_m^{1/2}(\log L_m)^2.
        \end{align*}
        Hence, for each fixed $\eta>0$,
        $$
        \sum_{m=1}^{\infty}
        \operatorname{meas}\{y\in J\mid T_m(y)>\eta L_m\}<\infty.
        $$
        Applying Lemma~\ref{lem-borel-cantelli} with $\eta=1/s$, and then taking the countable intersection over positive integers $s$, gives, for almost all $y\in J$,
        $$
        T_m(y)=o_y(L_m).
        $$
        Since
        $$
        \sum_{\substack{d\ge1\\ d\ \mathrm{squarefree}}}\frac1{d^2}
        =\prod_p\left(1+\frac1{p^2}\right)
        =\frac{\zeta(2)}{\zeta(4)},
        $$
        it follows from \eqref{eq-small-moduli-sharp} that
        $$
        \sum_{\substack{d\le L_m^B\\ d\ \mathrm{squarefree}}}\frac{M_{L_m,d}(y)}d
        \le (\zeta(2)/\zeta(4)+o_y(1))L_m
        $$
        for almost all $y\in J$.

        We next pass from $L_m$ to arbitrary $K$. If $L_m\le K<L_{m+1}$, monotonicity gives
        $$
        \begin{aligned}
            \sum_{\substack{d\le K^B\\ d\ \mathrm{squarefree}}}\frac{M_{K,d}(y)}d
            \le \sum_{\substack{d\le L_{m+1}^B\\ d\ \mathrm{squarefree}}}\frac{M_{L_{m+1},d}(y)}d
            \le (\zeta(2)/\zeta(4)+o_y(1))L_{m+1}.
        \end{aligned}
        $$
        Since $L_{m+1}\le(\lambda+o(1))K$, taking the countable intersection over $\lambda=1+1/s$, where $s=1,2,\ldots$, gives
        \begin{equation}\label{eq-small-ae-sharp}
            \sum_{\substack{d\le K^B\\ d\ \mathrm{squarefree}}}\frac{M_{K,d}(y)}d
            \le (\zeta(2)/\zeta(4)+o_y(1))K.
        \end{equation}

        For each fixed $d$, let $m_a$ be the number of indices $k\le K$ for which $\lfloor y^k\rfloor\equiv a\pmod d$. Then
        \begin{align*}
            \#\{1\le k<\ell\le K\mid \lfloor y^k\rfloor\equiv\lfloor y^\ell\rfloor\pmod d\}
            &=\sum_a\binom{m_a}{2}\\
            &\le \frac12\sum_a m_aM_{K,d}(y)\\
            &=\frac K2M_{K,d}(y).
        \end{align*}
        After summing with weight $1/d$ over squarefree $d\le K^B$, \eqref{eq-small-ae-sharp} gives
        \begin{equation}\label{eq-small-pair-sharp}
            \sum_{\substack{1\le k<\ell\le K\\ \lfloor y^k\rfloor\ne\lfloor y^\ell\rfloor}}
            \sum_{\substack{d\mid \lfloor y^\ell\rfloor-\lfloor y^k\rfloor\\d\le K^B\\ d\ \mathrm{squarefree}}}\frac1d
            \le \left(\frac{\zeta(2)}{2\zeta(4)}+o_y(1)\right)K^2.
        \end{equation}

        It remains to consider divisors $d>K^B$. For the same sequence $L_m$, define
        $$
        R_m(y)=\sum_{1\le k<\ell\le L_{m+1}}
        \sum_{\substack{L_m^B<d\le Y_2^\ell\\
                \lfloor y^k\rfloor\ne\lfloor y^\ell\rfloor\\
                \lfloor y^\ell\rfloor\equiv\lfloor y^k\rfloor\pmod d}}
        \frac1d.
        $$
        If the inner congruence holds and the two floors are different, then
        $$
        d\le \left|\lfloor y^\ell\rfloor-\lfloor y^k\rfloor\right|\le Y_2^\ell.
        $$
        By Lemma~\ref{lem-collision},
        \begin{align*}
            \int_J R_m(y)\,\mathrm{d}y
            &\ll \sum_{1\le k<\ell\le L_{m+1}}
            \sum_{L_m^B<d\le Y_2^\ell}
            \left(\frac1{d^2}+\frac1{d(\ell-k)Y_1^{\ell-1}}\right)\\
            &\ll \frac{L_{m+1}^2}{L_m^B}
            +\sum_{\substack{\ell\le L_{m+1}\\Y_2^\ell>L_m^B}}
            \frac{1+\log(Y_2^\ell/L_m^B)}{Y_1^{\ell-1}}
            \sum_{r=1}^{\ell-1}\frac1r.
        \end{align*}
        Put $\ell_0=B\log L_m/\log Y_2$. The second term is
        $$
        \ll_J \sum_{\ell>\ell_0}\frac{\ell\log \ell}{Y_1^\ell}
        \ll_J (\ell_0+1)^2Y_1^{-\ell_0}.
        $$
        By \eqref{eq-B-choice}, this is $O(L_m^{-3})$, while
        $$
        L_{m+1}^2/L_m^B\ll_\lambda L_m^{2-B}.
        $$
        Consequently, for every fixed $\eta>0$,
        $$
        \sum_{m=1}^{\infty}
        \operatorname{meas}\{y\in J\mid R_m(y)>\eta L_m^2\}<\infty.
        $$
        Applying Lemma~\ref{lem-borel-cantelli} with $\eta=1/s$, and then taking the countable intersection over positive integers $s$, gives
        $$
        R_m(y)=o_y(L_m^2)
        $$
        for almost all $y\in J$. If $L_m\le K<L_{m+1}$, every contribution with $k,\ell\le K$ and $d>K^B$ is included in $R_m(y)$. Hence
        \begin{equation}\label{eq-large-pair-negligible}
            \sum_{\substack{1\le k<\ell\le K\\ \lfloor y^k\rfloor\ne\lfloor y^\ell\rfloor}}
            \sum_{\substack{d\mid \lfloor y^\ell\rfloor-\lfloor y^k\rfloor\\d>K^B}}\frac1d
            =o_y(K^2).
        \end{equation}
        Finally,
        $$
        W(n)=\sum_{d\mid \operatorname{rad}(n)}\frac{1}{d}
        $$
        for every positive integer $n$. Combining \eqref{eq-small-pair-sharp} and \eqref{eq-large-pair-negligible} proves
        \begin{equation}\label{eq-pair-upper}
            \limsup_{K\to\infty}\frac1{K^2}
            \sum_{\substack{1\le k<\ell\le K\\
                    \lfloor y^k\rfloor\ne\lfloor y^\ell\rfloor}}
            W\!\left(\lfloor y^\ell\rfloor-\lfloor y^k\rfloor\right)
            \le \frac{\zeta(2)}{2\zeta(4)}.
        \end{equation}

        For the reverse inequality, intersect with the set of full measure on which Lemma~\ref{lem-residue-ud} holds. Fix $D\ge1$. For a squarefree integer $d\le D$ and $0\le r<d$, put
        $$
        N_r(K)=\#\{k\le K\mid \lfloor y^k\rfloor\equiv r\pmod d\}.
        $$
        By Lemma~\ref{lem-residue-ud}, for almost all $y\in J$,
        $$
        N_r(K)=\frac Kd+o_{y,d}(K).
        $$
        Hence, for every fixed squarefree $d$,
        $$
        C_d(K)=\sum_{r=0}^{d-1}\binom{N_r(K)}2
        $$
        satisfies
        $$
        C_d(K)=\frac12\sum_{r=0}^{d-1}N_r(K)^2-\frac K2
        $$
        and therefore
        $$
        C_d(K)=\frac{K^2}{2d}+o_{y,d}(K^2).
        $$
        Since $y^{k+1}-y^k=(y-1)y^k>1$ for all sufficiently large $k$, the sequence $\lfloor y^k\rfloor$ is eventually strictly increasing. Hence the number of pairs $k<\ell$ with $\lfloor y^k\rfloor=\lfloor y^\ell\rfloor$ is bounded in terms of $y$, and
        $$
        \#\{k<\ell\le K\mid \lfloor y^k\rfloor\ne\lfloor y^\ell\rfloor,\ d\mid \lfloor y^\ell\rfloor-\lfloor y^k\rfloor\}
        =\frac{K^2}{2d}+o_{y,d}(K^2).
        $$
        Keeping only the squarefree divisors not exceeding $D$ in the divisor expansion of $W$ gives
        $$
        \begin{aligned}
            &\sum_{\substack{k<\ell\le K\\ \lfloor y^k\rfloor\ne\lfloor y^\ell\rfloor}}
            W\!\left(\lfloor y^\ell\rfloor-\lfloor y^k\rfloor\right)\\
            &\qquad\qquad\ge
            \sum_{\substack{d\le D\\d\ \mathrm{squarefree}}}\frac1d
            \#\{k<\ell\le K\mid \lfloor y^k\rfloor\ne\lfloor y^\ell\rfloor,\ d\mid \lfloor y^\ell\rfloor-\lfloor y^k\rfloor\}.
        \end{aligned}
        $$
        Consequently
        $$
        \liminf_{K\to\infty}\frac1{K^2}
        \sum_{\substack{k<\ell\le K\\ \lfloor y^k\rfloor\ne\lfloor y^\ell\rfloor}}
        W\!\left(\lfloor y^\ell\rfloor-\lfloor y^k\rfloor\right)
        \ge
        \frac12\sum_{\substack{d\le D\\d\ \mathrm{squarefree}}}\frac1{d^2}.
        $$
        Letting $D$ tend to infinity gives
        $$
        \liminf_{K\to\infty}\frac1{K^2}
        \sum_{\substack{k<\ell\le K\\ \lfloor y^k\rfloor\ne\lfloor y^\ell\rfloor}}
        W\!\left(\lfloor y^\ell\rfloor-\lfloor y^k\rfloor\right)
        \ge \frac{\zeta(2)}{2\zeta(4)}.
        $$
        Together with \eqref{eq-pair-upper}, this gives \eqref{eq-pair-asymptotic}.
    \end{proof}

\end{document}
